% TOP_PROBLEM: {"id": 7900001, "title": "Extended States in the Anderson Model", "category_id": 16, "category": {"id": 16, "name": "physics", "display_name": "Mathematical Physics", "description": "Problems at the intersection of mathematics and physics.", "slug": "physics", "order_index": 16, "created_at": "2026-07-31T15:26:25.671Z"}, "status": "open"}
% FIRST_POSTED: 2026-09-27
% !TeX program = lualatex
% Standalone notebook; compile twice: lualatex 7900001.tex
% Original section preserved verbatim except for marked insertion blocks.
\documentclass[11pt,a4paper]{article}
\usepackage[margin=24mm,headheight=14pt]{geometry}
\usepackage{fontspec}
\setmainfont{Latin Modern Roman}
\setsansfont{Latin Modern Sans}
\setmonofont{Latin Modern Mono}
\usepackage{amsmath,amssymb,mathtools}
\usepackage{unicode-math}
\setmathfont{Latin Modern Math}
\usepackage{microtype}
\usepackage{enumitem,xurl,needspace}
\usepackage[dvipsnames]{xcolor}
\usepackage{hyperref}
\hypersetup{unicode=true,colorlinks=true,linkcolor=MidnightBlue,urlcolor=MidnightBlue,
 pdftitle={Research notebook: Problem 7900001},
 pdfauthor={Open Problems Math research notebook},bookmarksnumbered=true}
\usepackage{bookmark,titlesec,fancyhdr}
\titleformat{\section}{\Large\bfseries\raggedright}{\thesection}{0.7em}{}
\pagestyle{fancy}\fancyhf{}
\fancyhead[L]{\small\sffamily Open Problems Math\enspace|\enspace Research notebook}
\fancyhead[R]{\small\sffamily Problem 7900001}
\fancyfoot[L]{\small\sffamily 27 September 2026}
\fancyfoot[R]{\small\sffamily \thepage}
\setlength{\parindent}{0pt}
\setlength{\parskip}{5pt plus 1pt}
\setlength{\emergencystretch}{3em}
\setcounter{secnumdepth}{1}
\setcounter{tocdepth}{1}
\setlist[itemize]{leftmargin=3em,itemsep=5pt,topsep=4pt,parsep=0pt}
\allowdisplaybreaks
\newcommand{\N}{\mathbb{N}}
\newcommand{\Z}{\mathbb{Z}}
\newcommand{\Q}{\mathbb{Q}}
\newcommand{\R}{\mathbb{R}}
\newcommand{\C}{\mathbb{C}}
\newcommand{\F}{\mathbb{F}}
\newcommand{\A}{\mathbb{A}}
\newcommand{\PP}{\mathbb P}
\newcommand{\E}{\mathbb{E}}
\newcommand{\eps}{\varepsilon}
\newcommand{\dd}{\,\mathrm d}
\newcommand{\sref}[2]{\hyperlink{source:#1:#2}{[S#2]}}
\newcommand{\eref}[2]{\hyperlink{extra:#1:#2}{[E#2]}}
\newif\ifshowcatalogue
\showcataloguetrue
\newcommand{\cataloguescope}[1]{\ifshowcatalogue
 \paragraph{Exact scope in the supplied catalogue.}
 \begin{quote}\small #1\end{quote}\fi}
\numberwithin{equation}{section}
\begin{document}
\setcounter{section}{0}
% BEGIN PRESERVED SECTION
% ================================================================
% problemId: 7900001
% Canonical problem ID: 7900001
% exactTarget SHA-256: 7c69569d51e15ba2f9566f102269741fba36f1d2600bcceb9a4b69c220458fd5
\clearpage
\section*{\texorpdfstring{Extended States in the Anderson Model}{Extended States in the Anderson Model}}\label{7900001}
\subsection{Definitions and mathematical statement}
For independent identically distributed real random variables $V_x$, the standard lattice Anderson operator on $\ell^2(\Z^d)$ has the form
\[
 (H_\lambda\psi)(x)=-\sum_{|y-x|_1=1}\psi(y)+\lambda V_x\psi(x).
\]
For a vector $\psi$, its spectral measure is $\mu_\psi(B)=\langle\psi,1_B(H_\lambda)\psi\rangle$; an absolutely continuous spectral component has density with respect to Lebesgue measure. The central target is a nonzero absolutely continuous spectral subspace in a weak-disorder regime in $d\ge3$. The catalogue does not specify a probability law or exact energy interval. Those choices must be fixed in a full model-specific theorem. Its additional language about transport and a mobility edge is not, by definition alone, equivalent to existence of absolutely continuous spectrum.
\par\smallskip\textit{Formulation and concept references:} \sref{7900001}{1}.
\cataloguescope{Prove that the Anderson model, the discrete random Schroedinger operator on Z\textasciicircum{}d with d >= 3, has a regime of absolutely continuous spectrum, equivalently that a conducting phase and a mobility edge exist at weak disorder, so that the metal-insulator transition is established.}
\subsection{Short English statement}
Can weak random disorder in a lattice of dimension at least three be proved to permit an extended-state spectral regime?
% BEGIN RESEARCH 31
\subsection{Research attempt}

\paragraph{Current frontier.}
Kirsch distinguishes spectral type, dynamical localization and averaged
density-of-states information; these are different assertions.
\sref{7900001}{1}. Black--Drogin--Hern\'andez establish weak-disorder diffusion
for Gaussian disorder down to a stated imaginary-resolvent scale
$\eta\asymp\lambda^{2+\kappa_d}$, with $\kappa_d>0$, in their November 2025
version. It does not let $\eta\downarrow0$ at fixed positive $\lambda$.
\sref{7900001}{3}, \eref{7900001}{1}. Jak\v{s}i\'c--Last--Warzel's September 2026
manuscript proves a Bethe-lattice result and proposes a Poisson--Schur
criterion on $\Z^d$, explicitly leaving its weak-disorder lower bound
unproved there. \sref{7900001}{2}, \eref{7900001}{2}. Becker--Oltman's September
2026 manuscript concerns a hyperbolic-square-lattice product, not $\Z^d$
alone. \eref{7900001}{3}. Recent statements and their model scopes were checked;
they are not treated as independently audited solutions of this entry.

\paragraph{Attack route.}
\textbf{Editorial scope note.} The original does not prescribe a single-site
law and itself warns that absolutely continuous spectrum is not equivalent,
by definition, to transport or a mobility edge. To make a concrete attempt,
fix $d\ge3$ and independent $V_x$ uniformly distributed on $[-1,1]$, with
$\lambda>0$. This is a nondegenerate random model, not the trivial free
case. Results below concern its spectral component only; they do not
silently settle the additional transport language or all possible laws.

The selected route is to upgrade resolvent control to a quenched statement
about individual spectral measures. A uniform $L^2$ bound on their Poisson
transforms is stronger than needed, but it gives a clean sufficient estimate
whose limiting steps can be proved. The competing tree-cyclicity route
would instead need a lower bound on a matrix Schur complement; independence
of forward branches is absent on $\Z^d$. \eref{7900001}{2}, Section~4.
I investigate the first route and test whether averaging or small operator
norm perturbations could supply its missing estimate.

\paragraph{Attempt.}
Write
\[
 G_x(E+i\eta)=\langle\delta_x,(H_\lambda-E-i\eta)^{-1}\delta_x\rangle,
 \qquad u_{x,\eta}(E)=\operatorname{Im}G_x(E+i\eta).
\]
For the probability measure $\mu_x$ associated to $\delta_x$, the spectral
theorem gives
\begin{equation}\label{r31:poisson}
 u_{x,\eta}(E)=\int\frac{\eta}{(t-E)^2+\eta^2}\,d\mu_x(t)\ge0.
\end{equation}
The absence of $1/\pi$ in this definition will be retained in all limits.
Spectral measures and the resolvent framework are those of \sref{7900001}{1}.

\textbf{Lemma 31.1 (an all-resolution sufficient estimate).}
Let $I$ be an open interval. For fixed $\lambda>0$, suppose that for every
closed interval $J\Subset I$,
\begin{equation}\label{r31:missing}
 \sup_{0<\eta<1}\mathbb E\int_J u_{0,\eta}(E)^2\,dE<\infty.
\end{equation}
Then, almost surely, $H_\lambda$ has no singular spectral component on $I$.

\emph{Proof.} Fix $J$ and put $\eta_l=2^{-l}$. Fatou's lemma implies
\[
 \mathbb E\left[\liminf_l\|u_{0,\eta_l}\|_{L^2(J)}^2\right]
       \le\liminf_l\mathbb E\|u_{0,\eta_l}\|_{L^2(J)}^2<\infty.
\]
Thus for almost every realization there is a subsequence with bounded
$L^2(J)$ norm. It need not be a deterministic subsequence or a subsequence
chosen independently of the realization. Weak compactness of $L^2(J)$ gives
a further weakly convergent subsequence with limit $h\in L^2(J)$.
For $\phi\in C_c(J^\circ)$, Fubini and the Poisson approximate identity give
\[
 \lim_{\eta\downarrow0}\int\phi(E)u_{0,\eta}(E)\,dE
                 =\pi\int\phi(t)\,d\mu_0(t).
\]
Indeed, the inner Poisson integral converges uniformly to $\pi\phi(t)$
and is bounded by $\pi\|\phi\|_\infty$, so integration against the finite
measure is justified. Weak $L^2$ convergence identifies the same limit as
$\int\phi h$. Hence $\mu_0|_{J^\circ}=h(E)dE/\pi$.

Translation invariance of the law gives the same conclusion for each $x$.
Intersect the probability-one events for countably many $x\in\Z^d$ and a
countable exhaustion of $I$ by such intervals. All basis-vector spectral
measures are then absolutely continuous on $I$. The singular spectral
projection on $I$ annihilates every $\delta_x$, and thus, by density, the
whole Hilbert space. This proves the lemma without interchanging an
uncontrolled random pointwise limit and expectation.\hfill$\square$

\textbf{Lemma 31.2 (nonvanishing is available without a new randomness estimate).}
Choose a nonnegative $\phi\in C_c(I)$ such that
$a_0=\langle\delta_0,\phi(H_0)\delta_0\rangle>0$. There is $\lambda_1>0$
so that for every $0<\lambda<\lambda_1$ and every realization of bounded
$|V_x|\le1$,
\begin{equation}\label{r31:mass}
           \langle\delta_0,\phi(H_\lambda)\delta_0\rangle>a_0/2.
\end{equation}
\emph{Proof.} The operator norm of $H_\lambda-H_0$ is at most $\lambda$.
For $\lambda\le1$, both spectra lie in $[-2d-1,2d+1]$. Approximate $\phi$
uniformly on this interval by a polynomial $P$. For a fixed polynomial,
\[
 \|A^j-B^j\|\le j\max(\|A\|,\|B\|)^{j-1}\|A-B\|
\]
follows from the telescoping identity, with no commutativity required.
Consequently $\|P(H_\lambda)-P(H_0)\|\to0$ uniformly in the realization.
Choosing the approximation error first, and then $\lambda$, proves
\eqref{r31:mass}.\hfill$\square$

Such a $\phi$ exists for every nonempty open $I\subset(-2d,2d)$: Fourier
transform diagonalizes $H_0$ as multiplication by
$-2\sum_{j=1}^d\cos\theta_j$, and the preimage of a nonempty open energy
subinterval has positive torus measure. Thus proving
\eqref{r31:missing} for a nonempty $I$ and every sufficiently small positive
$\lambda$ would establish a nonzero absolutely continuous spectral subspace
for this specified model, indeed purity on $I$. The constants may depend
on $\lambda$; uniformity in \emph{resolution} $\eta$ at each fixed $\lambda$
is the indispensable demand.

\medskip
\textbf{Attempted shortcut: average density instead of squared density.}
This shortcut is false. Let $V$ be uniform on $[-1,1]$ and take the
one-dimensional operator with spectral measure $\mu=\delta_V$. Every
realization is purely atomic, while $\mathbb E\mu$ has density $1/2$ on
$[-1,1]$. Moreover
\[
 \mathbb E\frac{\eta}{(V-E)^2+\eta^2}\le\frac\pi2
 \quad\text{for all }E,\eta>0,
\]
yet
\begin{equation}\label{r31:atomic}
 \int_{\R}\left(\frac{\eta}{(V-E)^2+\eta^2}\right)^2dE
                   =\frac{\pi}{2\eta}.
\end{equation}
The last identity follows by $E-V=\eta t$ and
$\int_{\R}(1+t^2)^{-2}dt=\pi/2$. A bounded averaged Poisson density can
therefore conceal singular spectrum. This is not a counterexample to the
nearest-neighbor lattice conjecture; it disproves the proposed inference
from annealed to quenched spectral type.

\textbf{Attempted shortcut: a uniform Neumann expansion.}
Although $\|\lambda V\|\le\lambda$, the free resolvent norm is at most
$1/\eta$. Expanding
\[
 (H_\lambda-z)^{-1}
  =(H_0-z)^{-1}\bigl[I+\lambda V(H_0-z)^{-1}\bigr]^{-1}
\]
by its geometric series is guaranteed only when $\lambda/\eta<1$.
That regime excludes the required limit $\eta\downarrow0$ with fixed
$\lambda$. Averaging may cancel some terms, but a proof of a remainder
uniform in $\eta$ is an additional estimate, not supplied by the norm
bound. The controlled mesoscopic scale in \eref{7900001}{1} does not bridge
this gap, and its Gaussian law is not the bounded law chosen above.

\paragraph{Stress test.}
The order of limits cannot be changed unnoticed. Every finite-volume spectral
measure is atomic. Its positive squared Poisson transform has diagonal
contributions of order $\eta^{-1}$, just as in \eqref{r31:atomic}; therefore
one cannot prove the required all-resolution bound at each fixed finite
volume and then take volume to infinity. Nor does the small-$\lambda$
continuity in Lemma~31.2 preserve spectral type: it preserves only positive
mass for a fixed continuous test function. Lemma~31.1 supplies the separate
absolute-continuity step conditionally.

No numerical transport experiment is interpreted as a spectral proof here.
The companion script symbolically checks the Poisson-square integral.
The probability-one conclusion was derived for each fixed $\lambda$; no
uncountable intersection over all disorder strengths is asserted. Finally,
existence of an ac subspace, a quantitative conductivity statement, and a
sharp mobility edge remain separate goals, as the original editorial
formulation already notes.

\Needspace{8\baselineskip}
\paragraph{Outcome.}
\textbf{Outcome: Conditional reduction.}

For an explicitly fixed nondegenerate bounded single-site law, an
all-resolution second-moment resolvent estimate implies the desired
nonzero absolutely continuous spectral regime, with its nonvanishing and
almost-sure conclusions proved in detail. The estimate itself is not proved.
The averaging and norm-perturbation shortcuts are rejected by an exact model
calculation and a parameter-order obstruction. This is a sufficient route
for the spectral part of a specified model, not a completion of the
underspecified law/transport/mobility-edge package, and the criterion's
historical novelty is not asserted.

% END RESEARCH 31

\subsection{Sources}
\begin{itemize}\raggedright
\item[\textnormal{[S1]}]\hypertarget{source:7900001:1}{} W. Kirsch, "An Invitation to Random Schrodinger Operators", in Random Schrodinger Operators, Panoramas et Syntheses 25, Societe Mathematique de France, 2008; arXiv:0709.3707.
\url{https://arxiv.org/abs/0709.3707}.
{\small\textit{Catalogue formulation source.}}
\item[\textnormal{[S2]}]\hypertarget{source:7900001:2}{} A Note on Extended States on the Bethe Lattice.
\url{https://arxiv.org/html/2609.06987v1}.
{\small\textit{Catalogue research/status reference; not a proof endorsement.}}
\item[\textnormal{[S3]}]\hypertarget{source:7900001:3}{} Self-consistent equations and quantum diffusion for the Anderson model.
\url{https://arxiv.org/abs/2506.06468}.
{\small\textit{Catalogue research/status reference; not a proof endorsement.}}
% BEGIN ADDED REFERENCES 31
\item[\textnormal{[E1]}]\hypertarget{extra:7900001:1}{} Adam Black, Reuben
Drogin and Felipe Hern\'andez, \emph{Self-consistent equations and quantum
diffusion for the Anderson model}, arXiv:2506.06468v2,
6 November 2025, abstract and resolution-scale statement.
\url{https://arxiv.org/abs/2506.06468v2}.
{\small\textit{Gaussian disorder and fixed scale are retained; checked
27 September 2026.}}
\item[\textnormal{[E2]}]\hypertarget{extra:7900001:2}{} Vojkan Jak\v{s}i\'c,
Yoram Last and Simone Warzel, \emph{A Note on Extended States on the
Bethe Lattice}, arXiv:2609.06987v1, Sections~1 and 4,
Lemma~4.1 and Theorem~4.2.
\url{https://arxiv.org/html/2609.06987v1}.
{\small\textit{The lattice criterion is not its unproved weak-disorder
estimate. Consulted 27 September 2026.}}
\item[\textnormal{[E3]}]\hypertarget{extra:7900001:3}{} Simon Becker and
Izak Oltman, \emph{Disorder on the hyperbolic square lattice I:
Anderson delocalization and absolutely continuous spectrum},
arXiv:2609.20798v1, September 2026, abstract and graph hypotheses.
\url{https://arxiv.org/abs/2609.20798v1}.
{\small\textit{Recent manuscript on a different graph; checked
27 September 2026, not fully audited here.}}

% END ADDED REFERENCES 31
\end{itemize}

% END PRESERVED SECTION
\end{document}
